% \code (common): inline code on a pill that may break across lines, at its capitals, after :: and _, and at
% spaces, without a hyphen. Its text is literal, read as an argument (a frame without [fragile]) and read
% directly (a [fragile] frame, where % is literal too). The pills do not cover the ink of their glyphs (each word of
% the same red text is as wide in \code as in \texttt). It is also used in a section title (PDF bookmark) and with
% draftpics.
\documentclass[aspectratio=43,9pt]{beamer}
\usetheme{joby}
\usepackage{common}
\usepackage[draftpics]{draft-tools}
\title{Code}
\author{Test Author}
\begin{document}
\section{About \code{CSectionName}}
\begin{frame}{Code}
	Short \code{CMessage::Send()} and words \code{some_name and more}.

	\parbox{0.4\linewidth}{Long \code{CInputSignalManagementTelemetryBase::DoSomething()}}
\end{frame}
\begin{frame}{Literal}
	LIT \code{f(a){b}} \code{arr[i][j]} \code{user@host} \code{cost$5} \code{a&b} \code{x^2} \code{a~b}
	\code{\textbf{x}} \code{some\_name} \code{\a\b c}
\end{frame}
\begin{frame}[fragile]{Literal read directly}
	DIRECT \code{50% #define \x{y}}
\end{frame}
\begin{frame}{Ink}
	\colorlet{code text}{red}
	\code{thm mmm}

	\textcolor{red}{\texttt{thm mmm}}
\end{frame}
\end{document}
